\documentclass[10pt,letterpaper]{article}

% Font
\usepackage[T1]{fontenc}
\usepackage{mathpazo}

% Page layout
\usepackage{geometry}
\geometry{top=0.65in, bottom=0.75in, left=0.75in, right=0.75in}

% Spacing
\usepackage{setspace}
\setstretch{1.05}
\usepackage{microtype}

% Colors
\usepackage[x11names]{xcolor}

% Header/footer
\usepackage{fancyhdr}
\pagestyle{fancy}
\renewcommand{\headrule}{}
\fancyhf{}
\fancyfoot[C]{\footnotesize\textcolor{gray}{joros@cmu.edu}}
\fancyfoot[L]{\footnotesize\textcolor{gray}{sites.google.com/view/joseoros}}
\fancyfoot[R]{\footnotesize\textcolor{gray}{March 2026}}

% Section formatting
\usepackage{titlesec}
\titleformat*{\section}{\centering\LARGE\scshape}
\titleformat*{\subsection}{\bfseries\scshape\color{gray}}
\titlespacing*{\section}{0pt}{10pt}{4pt}
\titlespacing*{\subsection}{0pt}{7pt}{3pt}

% Lists
\usepackage{enumitem}
\setlist[itemize]{leftmargin=1.5em, label=\color{gray}{\textbullet}, nosep, topsep=2pt}

% Links
\usepackage[hyphens]{url}
\urlstyle{same}
\usepackage[hidelinks, hypertexnames=false]{hyperref}

% Custom commands
\newcommand{\name}[1]{\section*{#1}}
\newcommand{\sep}{{\color{gray}\,\textbullet\,\,}}
\newcommand{\hsep}{\noindent\textcolor{gray}{\rule{\textwidth}{0.5pt}}}
\newcommand{\block}[1]{\begin{samepage}\hsep\subsection*{#1}\end{samepage}}

\begin{document}

% Name
\name{Jose Eduardo Oros Chavarria}

\begin{center}
(412) 500-4169 \sep \href{mailto:joros@cmu.edu}{joros@cmu.edu} \sep \href{https://www.linkedin.com/in/joseoros}{linkedin.com/in/joseoros} \sep \href{https://sites.google.com/view/joseoros}{sites.google.com/view/joseoros}
\end{center}

% Summary
\block{Summary}

\noindent Data scientist with expertise in experimentation design, causal inference, and machine learning. Industry experience driving subscriber growth, retention, and product insights at Pandora and Spotify. PhD research on how generative AI systems and digital platforms reshape user behavior and creative markets.

% Skills
\block{Skills}

\noindent\textbf{Programming:} Python (pandas, NumPy, scikit-learn, statsmodels), R (tidyverse, fixest, ggplot2), SQL (PostgreSQL, BigQuery), Java \\
\textbf{Methods:} A/B Testing \& Experimentation, Causal Inference (DiD, IV, RDD, Synthetic Controls), GLMs, Survival Models, Forecasting, Machine Learning, NLP \\
\textbf{Tools \& Platforms:} AWS, BigQuery, Airflow, Tableau, Git \\
\textbf{Domain:} Product Analytics, User Behavior Analysis, Retention \& Growth, Platform Economics

% Professional Experience
\block{Professional Experience}

\begin{itemize}
\item \textbf{Carnegie Mellon University} \sep Doctoral Researcher \sep Pittsburgh, PA \hfill Aug 2019 -- present
  \begin{itemize}[label={}]
  \item Designed and executed a field experiment with 1,000+ participants to measure how LLM-assisted ideation affects the quality and originality of creative output.
  \item Built ML prediction models on large-scale social media data to forecast career outcomes for emerging musical artists.
  \item Led the establishment of a multi-year data partnership with Spotify's Creator Economics team; analyzed proprietary behavioral data to quantify how platform promotion drives on- and off-platform outcomes using causal inference (DiD, IV).
  \item Applied causal inference methods (DiD, IV, synthetic controls) to large observational datasets to identify platform policy effects.
  \item Presented findings to cross-functional audiences at Spotify, MIT, INFORMS, and NBER (10+ talks).
  \end{itemize}
\item \textbf{Spotify} \sep Economics Research Intern, Creator Economics \sep New York City, NY \hfill Summer 2022
  \begin{itemize}[label={}]
  \item Partnered with the Creator Economics product team to analyze artist performance and platform impact using proprietary data and causal inference methods.
  \item Scoped and launched an ongoing multi-year data-sharing partnership between Spotify and Carnegie Mellon, shaping the research agenda for platform creator economics.
  \end{itemize}
\item \textbf{Pandora Media} \sep Data Analyst, Subscriber Growth \& Retention \sep Oakland, CA \hfill Jan 2018 -- Aug 2019
  \begin{itemize}[label={}]
  \item Designed and analyzed A/B tests for subscriber acquisition and retention campaigns; built forecasting models to measure and optimize marketing ROI using R, Python, and SQL.
  \item Built Tableau dashboards to monitor subscriber KPIs for the Marketing Analytics team, enabling data-driven campaign decisions across channels.
  \item Developed and maintained data pipelines using Airflow on AWS to process digital advertising engagement data at scale.
  \end{itemize}
\end{itemize}

% Education
\block{Education}

\begin{itemize}
\item \textbf{Carnegie Mellon University}, H.\ John Heinz III College, Pittsburgh, PA
  \begin{itemize}[label={}]
  \item Ph.D., Information Systems and Management \hfill 2026 (expected)
  \item Daniel and Rise Nagin Dissertation Fellowship \hfill 2025
  \item Master of Information Systems Management \hfill Dec 2017
  \end{itemize}
\item \textbf{Instituto Tecnol\'ogico Aut\'onomo de Mexico (ITAM)}, Mexico City, Mexico \hfill 2016
  \begin{itemize}[label={}]
  \item Bachelor of Science, Telecommunications Engineering
  \end{itemize}
\end{itemize}

\newpage

% Research Projects
\block{Research Projects}

\begin{itemize}
\item \textbf{AI and Creativity} --- Designed a controlled experiment to measure how LLM-generated ideas affect the quality and originality of music compositions. Found that AI assistance can reduce creative diversity, with implications for AI-augmented product workflows.
\item \textbf{Platform Promotion Spillovers} --- Used causal inference on Spotify's proprietary data to estimate how editorial playlist placement affects artists' streaming, concert attendance, and social media growth.
\item \textbf{Predicting Career Outcomes} --- Built machine learning models using large-scale social media data to forecast which emerging artists will achieve commercial success.
\end{itemize}

% Recitation Instructor
\block{Recitation Instructor}

\begin{itemize}
\item \textbf{Advanced Business Analytics}, Carnegie Mellon University \hfill 2021 -- 2025
  \begin{itemize}[label={}]
  \item Taught generalized linear models, survival analysis, and predictive modeling to classes of 10--20 graduate students. Four consecutive years.
  \end{itemize}
\end{itemize}

% Selected Presentations
\block{Selected Presentations}

\noindent 10+ invited talks at MIT (BIG.AI, CODE@MIT), INFORMS, NBER, Spotify, and international forums on generative AI and platform economics (2023--2026).

% Languages
\block{Languages}

\noindent English, Spanish (Native)

\end{document}
